\clearpage
\renewcommand{\sectioncolor}{qbit}
\section{On QBIT --- scoping the ambition honestly}
\markboth{Scoping QBIT}{}

You built \bk{TensoQi101} and \bk{QBIT-HIS} to reach \emph{Quantum Biological Information Theory} and
attack three problems with it: mutation, docking, structure prediction. The libraries are excellent.
The scoping deserves to be precise, because getting it wrong costs years.

\begin{center}\small
\begin{tabular}{@{}>{\raggedright\arraybackslash}p{4.0cm}c>{\raggedright\arraybackslash}p{11.3cm}@{}}
\toprule
\rowcolor{qbit!13}
\textbf{\color{qbit}Problem} & \textbf{quantum?} & \textbf{\color{qbit}What it actually needs}\\
\midrule
\textbf{Protein structure prediction} & \no &
 Classical statistical mechanics and geometry. \textbf{There is no quantum mechanics in AlphaFold.}
 Your route here is PBoC2 $\to$ Bressloff $\to$ the transfer-operator material of Part IX.\\
\addlinespace[2pt]
\textbf{Protein--ligand docking} & \pt &
 Classical statistical mechanics for the free energy; quantum only to \emph{parameterise} $U(x)$.
 \textbf{That is precisely what your four DFT texts are for} --- and why GROMACS is on your machine.\\
\addlinespace[2pt]
\rowcolor{qbit!7}
\textbf{Mutation} & \yes &
 \textbf{Genuinely quantum, in one specific place}: proton tunnelling in base pairs (Löwdin 1963;
 Slocombe, Sacchi \& Al-Khalili 2022). But most spontaneous mutation is classical --- replication
 error, deamination, oxidative damage.\\
\bottomrule
\end{tabular}
\end{center}

\begin{keyidea}[title={The honest shape of QBIT}]
\begin{center}\itshape\large
Quantum information theory is the right language for\\one third of one of your three problems.
\end{center}
\vspace{3pt}
That is not a demotion of the ambition --- it is a sharpening of it. Two consequences:
\begin{enumerate}[label=\textbf{\arabic*.}]
\item \textbf{The information-theoretic angle that generalises across all three is Shannon and
 Kolmogorov, not qubits.} Entropy, mutual information, large deviations and rate functions apply to
 gene expression, ligand binding and conformational ensembles alike. Nielsen \& Chuang \textbf{ch.~11}
 gives you the classical half before the quantum half.
\item \textbf{Where quantum genuinely enters --- the tunnelling substrate of mutation --- your real
 obstacle is Part IX's Tier 3:} non-Markovian open-system dynamics at strong coupling, in a
 structured bath. \textbf{N\&C ch.~8} is the entry; \textbf{Bru \& de Siqueira Pedra (2020)} is the
 rigorous destination; HEOM is in neither and would have to be acquired.
\end{enumerate}
\end{keyidea}

\begin{critbox}[title={Two books in QBIT-HIS to read with your guard up}]
\textbf{Khrennikov, \textit{Open Quantum Systems in Biology, Cognitive and Social Sciences} (2023).}
The open-systems mathematics is standard and useful. The programme of applying quantum formalism to
cognition and social science is \textbf{idiosyncratic and not accepted}. Take the formalism; leave
the interpretation.\\[5pt]
\textbf{\textit{Towards Quantum Integrated Information Theory} (1806.01421).} Speculative, and built
on IIT, which is itself contested. Not a foundation to build on.\\[5pt]
By contrast \textbf{Marais et al., \textit{The future of quantum biology} (2018)} --- also in
\bk{QBIT-HIS} --- is a legitimate review and the right place to calibrate what is established versus
what is hoped for. Read it before the other two.
\end{critbox}

\begin{plain}
The most valuable thing in your QBIT folder is not the quantum material. It is \textit{Physical
Biology of the Cell} and Bressloff --- statistical mechanics applied to real cells, with numbers.
Those two books connect to all three of your target problems. The quantum books connect to one
corner of one of them, and that corner happens to be the hardest open problem in the field.
\end{plain}
